\section{The wall turns on at a single operation}
\label{sec:depth}

Section~\ref{sec:wall} contrasted readable and computed errors as two categories. This section varies the distance between them continuously and finds that, for the strongest model, the transition is a step.

We vary the \emph{verification depth} of a computed error: the number of operations that must be redone to recover the true value from re-readable inputs. Depth zero is a readable error, whose truth is stated verbatim. Depth one places one operation between the planted value and its inputs, depth two places two, and so on. We hold the trace length and the number of tokens between the planted line and its nearest stated input fixed across depths, so depth measures the amount of computation a check requires and not the length of the trace or the distance the model must look back.

For the frontier model the result is a step function. Absorption is 0.01 at depth zero and 1.00 at depth one, and it stays at 1.00 through depths two, three, and five. A single operation between the model and the truth is enough to move it from catching essentially every error to catching none, and adding further operations changes nothing. The failure is not graded effort that decays as checking grows harder; it is a categorical refusal to recompute that the first operation already triggers in full.

The open-weight models reach the same ceiling by a shallower path. Llama-3.1-8B rises from 0.60 at depth zero to 0.99 at depth one, a sharp jump with a leakier readable baseline. Qwen2.5-7B ramps more gradually, from 0.71 at depth zero to 0.78 at depth one to 0.94 by depth two. The shape of the onset therefore depends on the model, from a clean step at the frontier to a short ramp for a mid-size model, but in every case absorption is at its ceiling by depth two and never falls again at greater depth. The height of the depth-zero starting point recapitulates Section~\ref{sec:wall}: the frontier model begins near zero because it catches readable errors, and the weaker models begin high because they do not.

Figure~\ref{fig:depth} plots absorption against verification depth for the three models, with the frontier step function beside the open-weight ramps.
